% Pista ana-23j-q1 · Anàlisi
% Procedència: PAU Catalunya 2023 · Convocatòria de juny · Sèrie 1 · Qüestió 1
\documentclass[11pt,a4paper]{article}
\usepackage{pista}
\pistainfo{Catalunya 2023}{Convocatòria de juny}{Sèrie 1}

\begin{document}

\section*{\textcolor{blauPista}{Qüestió 1}}

\noindent Calculeu els coeficients $a$, $b$, $c$ i $d$ de la funció $f(x) = ax^{3} + bx^{2} + cx + d$ si sabem que l'equació de la recta tangent a la gràfica de la funció $f$ en el punt d'inflexió $(1, 0)$ és $y = -3x + 3$ i que la funció té un extrem relatiu en el punt de la gràfica d'abscissa $x = 0$. \hspace*{\fill}\textnormal{[2,5~punts]}

\pista{Cada condició de l'enunciat es tradueix en una equació. Si les comptes bé, en són quatre:
\begin{itemize}
\item «$(1, 0)$ és un punt d'inflexió» amaga \emph{dues} condicions: $(1, 0)$ pertany a la gràfica, i $f''(1) = 0$;
\item la recta tangent en $(1, 0)$ té pendent $-3$: una condició sobre $f'(1)$;
\item «extrem relatiu en $x = 0$»: una condició sobre $f'(0)$.
\end{itemize}
Obtindràs un sistema de quatre equacions amb les incògnites $a$, $b$, $c$ i $d$.}

\end{document}
